\subsection*{Partie A}

On considère la fonction $f$ définie sur $\mathbb{R}$ par
\[f(x)=\left(\dfrac12 x + 1\right) \e^{-x}.\]

On note \mbox{$(\mathcal{C}_f)$} sa courbe représentative dans un repère orthonormé.

On admet que la fonction $f$ est dérivable sur $\mathbb{R}$ et on désigne par $f'$ sa fonction dérivée.

\begin{enumerate}
	\item Démontrer que, pour tout réel $x$,
	\[f'(x)=\left(-\dfrac12 x -\dfrac12\right) \e^{-x}.\]

	\item Étudier les variations de la fonction $f$ sur $\mathbb{R}$.

	\item L'affirmation \og{}Pour tout réel $x$, \mbox{$f(x)\leqslant 2$}\fg{} est-elle vraie ou fausse ?

	Justifier.
	\item Déterminer l'équation réduite de la tangente à \mbox{$(\mathcal{C}_f)$} au point d'abscisse 0.
\end{enumerate}
	\begin{center}
		\fcolorbox{black}{white}{
	\begin{minipage}[t]{0.8\linewidth}\vspace{0cm}
		\textsc{Aide aux calculs}
		\medskip

	\begin{tblr}{X[c]X[c]X[c]}
	\mbox{$\e^{1}\approx\np{2.7}$}&
	\mbox{$\e^{-1}\approx\np{0.4}$}&
	\mbox{$\e^{2}\approx\np{7.4}$}
	\end{tblr}
	\end{minipage}
}
	\end{center}

\subsection*{Partie B}

Soit $g$ la fonction polynôme du second degré définie par
\[g(x)=\np{0.5}x^2+bx+c\]
où $b$ et $c$ sont deux réels.

On crée la courbe représentée ci-dessous, constituée des deux éléments suivants :
\begin{itemize}[label=$\bullet$]
	\item Une partie de la courbe \mbox{$(\mathcal{C}_f)$} représentative de la fonction $f$ définie dans la partie A sur l'intervalle \mbox{$[-2;0]$}.

	\item Une partie de la courbe \mbox{$(\mathcal{C}_g)$} représentative de la fonction $g$  sur l'intervalle \mbox{$[0;3]$}.
\end{itemize}

\begin{center}
		\begin{tikzpicture}[xscale=1.3,yscale=1.3,rotate=0]
	\clip(-3.2,-0.6)rectangle(3.9,4.3);
	\draw[blue!30,thick](-3.2,-0.6)grid(3.9,4.3);
	\draw[line width=0.3mm,->, >=triangle 45](-3.2,0)--(3.9,0);
	\draw[line width=0.3mm,->, >=triangle 45](0,-0.6)--(0,4.3);
	\foreach \i in {-3,-2,-1,1,2,3,4}{\draw (\i,2pt)--(\i,-2pt) node[below ]{\np{\i}} (2pt,\i)--(-2pt,\i)node[left]{\np{\i}};}
	\draw(-2pt,-2pt) node[below left]{\np{0}};
	\draw[red](-1.5,1.5)node{\mbox{$(\mathcal{C}_f)$}};
	\draw[purple](2.6,2.2)node{\mbox{$(\mathcal{C}_g)$}};
	\node[draw,cross out, thick,minimum width=2pt, minimum height = 2pt] at (0,1){}; \draw(0,1)node[above right=1.5pt]{M};
	\draw[red, line width=0.3mm]plot[smooth,domain=-2:0]	(\x,{(\x/2+1)*exp(-\x)});
	\draw[purple, line width=0.3mm]plot[smooth,domain=0:3](\x,{\x*\x/2-\x/2+1});
\end{tikzpicture}

\end{center}
\vspace{5mm}
Le but de cette partie est de déterminer les valeurs des réels $b$ et $c$.
\begin{enumerate}
	\item Le point \mbox{M~$(0;1)$} appartient aux deux courbes  \mbox{$(\mathcal{C}_f)$} et  \mbox{$(\mathcal{C}_g)$}.

	Démontrer que \mbox{$c=1$}.
	\item
 Au point  \mbox{M~$(0;1)$} les courbes  \mbox{$(\mathcal{C}_f)$} et  \mbox{$(\mathcal{C}_g)$} admettent la même tangente.

 Déterminer la valeur de $b$.
\end{enumerate}
